% Canonical manuscript source.
\documentclass[11pt]{article}

\usepackage[T1]{fontenc}
\usepackage{lmodern}
\usepackage{amsmath,amssymb,amsthm}
\usepackage{booktabs}
\usepackage[margin=1in]{geometry}
\usepackage[hidelinks]{hyperref}

\newtheorem{theorem}{Theorem}
\newtheorem{lemma}[theorem]{Lemma}
\newtheorem{proposition}[theorem]{Proposition}
\theoremstyle{remark}
\newtheorem{remark}[theorem]{Remark}

\newcommand{\HH}{\mathbb H}
\newcommand{\OO}{\mathcal O}
\newcommand{\PP}{\mathcal P}
\newcommand{\QQ}{\mathcal Q}
\newcommand{\inner}[2]{\left\langle #1,#2\right\rangle_J}

\title{An Exact Classification of Ideal Right-Angled\\
$\mathrm{SW2}$-Polyhedra Through 84 Chambers}
\author{Xiaojian Zeng\\
\small PBC School of Finance, Tsinghua University, Beijing 100190, China\\
\small \texttt{xiaojianzeng@foxmail.com}}
\date{September 8, 2026}

\begin{document}
\maketitle

\begin{abstract}
Let $\PP=\mathrm{SW2}$ be the Scharlau--Walhorn Coxeter polyhedron with
the face numbering and reflective chamber tiling used by Allcock, Devlin,
Felikson, Kontorovich, and Whitehead.  We prove, by an exact exhaustive
classification, that every convex finite-volume ideal right-angled
$\PP$-polyhedron containing at most $84$ $\PP$-chambers contains exactly
$48$ chambers and is isometric to the ideal right-angled octahedron.
The classification covers both canonical branches of the decomposition in
their Proposition~7.10, including its non-reflective mixed-block branch.
All geometric predicates use integer Lorentz arithmetic, and completeness
follows from finite chamber- and block-expansion lemmas.  This is only a
bounded result: no argument controls polyhedra with more than $84$ chambers,
and Conjecture~7.11 of the cited work remains unresolved.
\end{abstract}

\section{Statement and relation to the unbounded problem}

Throughout, a hyperbolic polyhedron is convex and of finite volume.  Put
$\PP=\mathrm{SW2}$, with the face numbering of
\cite[Appendix~B]{allcock2026arithmetic}.  Following
\cite[Definition~2.1]{allcock2026arithmetic}, a $\PP$-polyhedron is a
convex finite union of $\PP$-chambers in the fixed reflective tiling of
$\HH^3$.  This is stronger than merely admitting an arbitrary tiling by
isometric copies of $\PP$.

Allcock et al. prove that the ideal right-angled octahedron $\OO$ is a
$\PP$-polyhedron and that every ideal right-angled $\PP$-polyhedron is
tiled by the right-angled triangular bipyramid $\tfrac14\OO$
\cite[Proposition~7.10]{allcock2026arithmetic}.  They conjecture that every
such polyhedron is tiled by full copies of $\OO$
\cite[Conjecture~7.11]{allcock2026arithmetic}.  We establish the following
finite subsidiary result.

\begin{theorem}\label{thm:main}
Let $\QQ$ be a convex finite-volume ideal right-angled $\PP$-polyhedron
consisting of at most $84$ $\PP$-chambers.  Then $\QQ$ consists of exactly
$48$ $\PP$-chambers and is isometric to $\OO$.  In particular, it is tiled
by one copy of $\OO$.
\end{theorem}

The restriction in Theorem~\ref{thm:main} is essential to the present
argument.  We neither bound the size of a minimal counterexample to the
unbounded conjecture nor classify abstract gluings of quarter-octahedra.
Whitehead's combinatorial characterization of abstract ideal right-angled
polyhedra \cite{whitehead2026note} gives a useful realization screen, but it
does not by itself certify convexity inside the fixed $\mathrm{SW2}$ chamber
tiling.  Accordingly, our final tests use the actual embedded Lorentz
half-spaces.

The proof is computer-assisted but exact.  Its organization is comparable
to exact finite-library results in discrete geometry that combine
symmetry-reduced exhaustive enumeration with integer predicates while
separating the bounded theorem from the ambient open problem
\cite{zeng2026finite}.  Here the additional task is to prove that the two
canonical block models from Proposition~7.10 exhaust the embedded
$\mathrm{SW2}$ geometry.

\section{An integral Lorentz model}

Use the bilinear form
\[
  \inner uv=-u_0v_0+u_1v_1+u_2v_2+u_3v_3
\]
on $\mathbb R^4$.  A spacelike vector $n$ represents the oriented
half-space $\inner nX\leq 0$ in the hyperboloid, or equivalently Klein,
model.  The following four integral normals represent the faces of the base
$\mathrm{SW2}$ chamber:
\begin{equation}\label{eq:normals}
\begin{aligned}
 n_1&=(0,-1,1,0), & n_2&=(1,1,1,1),\\
 n_3&=(0,0,0,-1),& n_4&=(0,0,-1,1).
\end{aligned}
\end{equation}
The reflection in $n$ is
\begin{equation}\label{eq:reflection}
 R_nX=X-2\frac{\inner nX}{\inner nn}n.
\end{equation}
Every matrix obtained from \eqref{eq:reflection} for the normals in
\eqref{eq:normals} is integral and preserves the Lorentz form.  Direct
inner-product calculation gives the Coxeter orders
\begin{equation}\label{eq:coxeter}
 (m_{ij})=
 \begin{pmatrix}
 1&2&2&3\\
 2&1&4&2\\
 2&4&1&4\\
 3&2&4&1
 \end{pmatrix}.
\end{equation}

The subgroup generated by $R_{n_1},R_{n_3},R_{n_4}$ has $48$ elements.
Closing the base chamber under this finite group and merging coplanar
boundary patches produces the eight oriented supporting planes
\begin{equation}\label{eq:standard-oct}
  (1,\varepsilon_1,\varepsilon_2,\varepsilon_3),
  \qquad \varepsilon_i\in\{-1,1\}.
\end{equation}
Their independent triple intersections give six projective null vectors,
and adjacent supporting normals are Lorentz-orthogonal.  Thus the union is
the standard ideal right-angled octahedron $\OO$.

The same calculation reconstructs both $12$-chamber copies of
$\tfrac14\OO$ used in Proposition~7.10.  Each has six facets, nine
right-angled edges, three ideal vertices, and two finite vertices.  Their
normalized facet Gram matrices agree after a permutation of the facets, so
the two quarters are isometric even though their $\mathrm{SW2}$ chamber
decompositions differ.

\section{Reduction to two finite censuses}

Lemma~7.9 of \cite{allcock2026arithmetic} gives a dichotomy for an ideal
right-angled $\PP$-polyhedron $\QQ$: either every boundary face is tiled by
face~3 of $\PP$, or no boundary face is.  Proposition~7.10 of the same work
then gives canonical decompositions as follows.

\begin{enumerate}
\item In the all-face-3 branch, $\QQ$ is reflectively tiled by
  $12$-chamber copies of $\tfrac14\OO$.
\item In the no-face-3 branch, $\QQ$ is canonically partitioned into
  $12$-, $24$-, and $48$-chamber copies of
  $\tfrac14\OO$, $\tfrac12\OO$, and $\OO$, respectively.  These blocks
  need not meet face-to-face as copies of the larger polyhedra, although
  their constituent $\PP$-chambers belong to the fixed reflective tiling.
\end{enumerate}

Call $12$ $\PP$-chambers one \emph{quarter unit}.  The hypothesis of
Theorem~\ref{thm:main} permits at most seven units.  The preceding dichotomy
is exhaustive, so it remains to classify the two branches through seven
units.

\subsection{The reflective-quarter branch}

Let $T=\tfrac14\OO$ with its right-angled reflection tiling.  We use the
following standard chamber-complex observation in a finite form.

\begin{lemma}[rooted interval expansion]\label{lem:interval}
Every connected convex union of at most $k$ chambers in a Coxeter complex,
after translating one chamber to the base chamber, is obtained by repeatedly
adjoining an adjacent chamber and closing under all chambers on all minimal
galleries.  It is contained in the Cayley ball of radius $k-1$.
\end{lemma}

\begin{proof}
Connectedness gives the radius bound.  An intersection of roots is closed
under every minimal gallery between two of its chambers: such a gallery
crosses each separating wall at most once.  Conversely, interval closure is
the chamber form of root convexity.  If a convex target has not yet been
reached, choose a target chamber adjacent to the current connected subset.
The interval closure after adjoining it stays inside every root containing
the target.  Induction reaches the target.
\end{proof}

For $k=7$, the complete radius-six ball contains $2060$ $T$-chambers.
Applying Lemma~\ref{lem:interval} with exact reflection matrices gives the
first row of Table~\ref{tab:counts}.  The three ideal right-angled outputs at
size four are rooted occurrences of the same octahedron; all other outputs
have an exactly detected finite boundary vertex.

\begin{table}[ht]
\centering
\caption{Exact census counts by quarter unit.  Reflective counts retain
rooted symmetry duplicates; mixed counts are translation classes.}
\label{tab:counts}
\begin{tabular}{lrrrrrrr}
\toprule
Quarter units &1&2&3&4&5&6&7\\
\midrule
Reflective rooted convex unions &1&6&18&57&150&450&1218\\
Reflective ideal right-angled &0&0&0&3&0&0&0\\
Mixed interface-compatible classes &1&2&4&17&52&235&1054\\
Mixed convex classes &1&2&2&10&11&35&69\\
Mixed ideal right-angled &0&0&0&2&0&0&0\\
\bottomrule
\end{tabular}
\end{table}

\subsection{The mixed canonical-block branch}

Every $\PP$-chamber has a unique image of the finite $1$--$3$--$4$ vertex.
Call it the chamber's $B_3$ center.  The canonical blocks in the no-face-3
branch partition the chambers by this center.

\begin{lemma}[local block completeness]\label{lem:local}
At a fixed $B_3$ center there are exactly $25$ oriented canonical local
blocks: twelve quarters, twelve halves, and one full octahedron.  Their
oriented cut-plane descriptors are pairwise distinct.  Exactly three
quarters, six halves, and one full octahedron contain any specified chamber
at the center.
\end{lemma}

\begin{proof}
The $48$ chambers around the center form the octahedron in
\eqref{eq:standard-oct}.  The orbit of mirrors of types $1$ and $4$ through
the center consists of six unoriented planes, hence twelve oriented
half-spaces.  Intersecting the central octahedron with each oriented
half-space gives twelve distinct $24$-chamber halves.  Testing every
orthogonal pair of their primitive integral normals and retaining the
$12$-chamber sectors gives twelve distinct quarters.  Adding the uncut star
gives $12+12+1=25$.  Recording the oriented cut normals proves injectivity.
Incidence counting, checked directly on the $48$ integral chamber matrices,
gives the stated $3$, $6$, and $1$ choices through a specified chamber.
\end{proof}

Only type-2 chamber adjacencies change the $B_3$ center.  A potentially
dangerous pruning in a block enumeration is to insist that adjacent blocks
agree at every represented type-2 interface.  The following exact sign
calculation proves that this loses no convex assembly.

\begin{lemma}[interface obstruction]\label{lem:interface}
If both $B_3$ centers adjacent across a type-2 chamber facet occur in a
canonical assembly but that interface chamber is present on only one side,
the assembly is not convex.
\end{lemma}

\begin{proof}
At the base center $e_0=(1,0,0,0)$, equation~\eqref{eq:reflection} gives
\[
 R_{n_2}e_0=(2,1,1,1),\qquad
 \inner{n_2}{e_0}=-1,\qquad
 \inner{n_2}{R_{n_2}e_0}=1.
\]
Thus the neighboring center lies strictly beyond the oriented type-2
supporting plane.  Lorentz equivariance gives the same pair of signs at
every translated interface.  An unmatched interface is exposed, while the
assembly contains the neighboring center on its forbidden side, contradicting
convexity.
\end{proof}

The block adjacency graph of a connected chamber union is connected: moves
within one center have types $1$, $3$, or $4$, while type $2$ is exactly the
move between centers.  Starting with one representative of each block size,
the enumeration therefore adjoins every neighboring center and each of the
ten local choices from Lemma~\ref{lem:local} that contains the required
interface chamber.  Lemma~\ref{lem:interface} permits immediate rejection
of mismatches.  A state is the finite set of centers together with the
oriented cut descriptor at each center.  Lemma~\ref{lem:local} makes this
encoding injective.  Normalizing the state at each of its constituent
chambers gives a complete quotient by left translation in the
$\mathrm{SW2}$ Coxeter group.  Induction on the connected block graph proves
that every canonical mixed assembly of at most seven units is enumerated.

The compatible- and convex-state counts appear in the final three rows of
Table~\ref{tab:counts}.  There are exactly two ideal right-angled survivors,
both at four units.  Their canonical block-size patterns are $[4]$ and
$[1,1,1,1]$.

\subsection{Exact geometric filters and identification}

For completeness, we describe the geometric predicate applied in both
branches.  Exposed chamber facets lying in the same mirror are merged and
represented by a primitive oriented integral normal.  A finite chamber union
is convex precisely when every one of its chambers lies in every resulting
inward half-space.  Necessity is immediate.  For sufficiency, if a point of
the half-space intersection were outside the chamber union, a generic segment
from a chamber interior point to it would first leave the union across an
exposed facet and enter the forbidden side of that facet.

Every boundary vertex is recovered by taking the projective intersection of
each linearly independent triple of supporting planes and retaining the
intersections satisfying all half-spaces.  Its Lorentz norm is negative for
a finite vertex and zero for an ideal vertex.  Two boundary facets form an
edge when they contain two recovered vertices, and that edge is right-angled
exactly when the two facet normals have Lorentz product zero.  Thus all
acceptance and rejection predicates are integer equalities and inequalities;
there is no floating-point tolerance.

For each of the two mixed survivors, after ordering its eight boundary
normals, the full Gram matrix agrees with that of the eight normals in
\eqref{eq:standard-oct}.  Four listed normals have coordinate determinant
$-8$, so they span Lorentz four-space.  Equality of the full Gram matrices
therefore extends the correspondence of normals to a Lorentz isometry,
including the orientations of all half-spaces.  Both survivors are isometric
to $\OO$.  The reflective survivors are identified in the same way (or by
the uniqueness part of the ideal right-angled realization theorem).  This
proves Theorem~\ref{thm:main}.

\section{Reproducibility and limitations}

The two ancillary verifiers use only the Python standard library and exact
integer arithmetic.  From the directory containing the evidence files, the
stored machine-readable outputs are reproduced by
\begin{verbatim}
python3 evidence/sw2_quarter_census.py --output /tmp/quarter.json
PYTHONPATH=evidence python3 evidence/sw2_mixed_census.py --output /tmp/mixed.json
\end{verbatim}
The first program reconstructs the Lorentz source model and performs the
reflective-quarter census.  The second imports that exact model, reconstructs
all local mixed blocks, verifies Lemmas~\ref{lem:local} and
\ref{lem:interface}, and performs the mixed census.  The stored JSON records
the complete count tables, local completeness checks, survivor half-spaces,
and Gram certificates.  Independent reruns during mathematical review
reproduced both outputs byte for byte.

This computation proves no assertion beyond seven quarter units.  In
particular, failure to find a larger counterexample is not evidence of a
global size bound, and the result does not upgrade the general tiling in
Conjecture~7.11 to a reflective or face-to-face tiling.

\paragraph{Computational and AI assistance.}
OpenAI Codex assisted with the literature search, the design and writing of
the finite enumeration programs, adversarial organization of the exact
checks, and initial manuscript preparation.  Language-model output is not
used as mathematical evidence: the claimed finite classification rests on
the reductions above and the released deterministic exact verifiers.  No
proof assistant was used.

\bibliographystyle{plain}
\bibliography{references}

\end{document}
